
Lorsque $X$ est de caractéristique nulle, i.e. son faisceau structural est un faisceau de $Q$-algébres, le problème se résoud simplement en prenant
\begin{equation}\tag{1.2}
 \Lambda^i L_\bullet=\bigl(\overset{i}{\otimes}{}^{\mathbb L}L_\bullet\bigr)^{\mathrm{alt}},
\end{equation}
où l'on considère la puissance tensorielle $i$.ème (au sens des catégories dérivées) de $L_\bullet$ comme étant un objet de $D^-(\underline{O}_i)$, $\underline{O}_i$ étant l'algèbre du groupe sysmétrique $\mathfrak S_i$ à coéfficients dans $\underline{O}_X$, et où l'exposant alt désigne le passage à la ``partie alternée sous l'action de $\mathfrak S_i$'' (opération qui a un sens grâce à l'hypothèse de caractéristique nulle). Une formule de ce type était utilisée déjà dans [8] (cité par la suite [RRR]), pour la première démonstration du théorème de Riemann--Roch.

Sans hypothèse de caractéristique, lorsque $L_\bullet$ est nul en degrés $>0$, alors par Dold--Puppe [6] $L_\bullet$ est canoniquement homotope au complexe associé à un Module semi-simplicial $L'_\bullet$ canoniquement associé à $L_\bullet$, et un candidat assez naturel pour $\Lambda^iL_\bullet$ (pour $L_\bullet$ plat, disons) serait le complexe associé au Module semi-simplicial $\Lambda^iL'_\bullet$ (dont la composante de degré $n$ est $\Lambda^i L^{\prime}_n$). Mais pour étendre la définition à des complexes de $D^-(X)$ qui ne sont plus supposés nuls en degrés $>0$, il faudrait étudier le comportement de l'opération précédente par rapport à l'opération de translation $L_\bullet\mapsto L_\bullet[1]$ (ou \uline{suspension}, dans la terminologie de [6]); malheureusement, il ne semble pas qu'on ait un isomorphisme $\Lambda^i(L_\bullet[1])\simeq(\Lambda^i L_\bullet)[i]$ qui permettrait de faire de façon triviale l'extension désirée.

\section*{2. La formule de Riemann--Roch sans hypothèses projectives}

\textbf{2.1.} Une fois en possession d'une $\lambda$-structure naturelle sur les $K^\bullet(X)$, il est possible de formuler le théorème de Riemann--Roch, sous la forme envisagée dans Exp. VIII, pour n'importe quel morphisme $f:X\to Y$ propre et d'intersection complète, (et même, plus généralement, dans le cas d'un schéma relatif $X$ sur un topos localement annelé $Y$, $X$ étant encore soumis aux conditions relatives correspondantes). La classe de Todd de $f$ se définit grâce à la classe dans $K^\bullet(X)$ du complexe cotangent relatif $T_f$ pour $f$, construit dans Exp. 0 4.4, lequel complexe est ici parfait grâce à l'hypothèse locale faite sur $f$. Il semble clair que la démonstration de la formule de Riemann--Roch dans ce cas général demandera l'introduction d'idées essentiellement nouvelles, et il est possible que ces mêmes idées pourraient donner la clef 
